% ============================================================
%  C: Como Programar -- 6ª Edição | Deitel & Deitel
%  Capítulo 5 -- Funções em C
%  EXERCÍCIOS (5.8 -- 5.48)
%  Abordagem incremental: Caps. 1 + 2 + 3 + 4 + 5
% ============================================================
\documentclass[12pt,a4paper]{article}
\usepackage[utf8]{inputenc}
\usepackage[T1]{fontenc}
\usepackage[brazil]{babel}
\usepackage{geometry}
\geometry{top=2.5cm,bottom=2.5cm,left=3cm,right=3cm}
\usepackage{parskip}\setlength{\parskip}{6pt}\setlength{\parindent}{0pt}
\usepackage{xcolor,tcolorbox,enumitem,listings,fancyhdr,titlesec,hyperref,booktabs,array,amsmath}
\tcbuselibrary{skins,breakable}

\definecolor{deitelblue}{RGB}{0,70,127}
\definecolor{deitelgray}{RGB}{240,242,245}
\definecolor{deitelgreen}{RGB}{0,110,60}
\definecolor{deitellightblue}{RGB}{220,235,250}
\definecolor{deitelred}{RGB}{160,20,20}
\definecolor{deitellorange}{RGB}{200,90,0}

\newtcolorbox{boaprática}[1][]{colback=deitellightblue,colframe=deitelblue,
  fonttitle=\bfseries\small,title={Boa prática de programação},breakable,#1}
\newtcolorbox{erroco}[1][]{colback=red!5,colframe=deitelred,
  fonttitle=\bfseries\small,title={Erro comum de programação},breakable,#1}
\newtcolorbox{respostabox}[1][]{colback=deitelgray,colframe=deitelblue!60,
  fonttitle=\bfseries\small\color{deitelblue},title={Resposta detalhada},breakable,#1}
\newtcolorbox{enunciadoBox}[1][]{colback=white,colframe=deitelblue,
  fonttitle=\bfseries,title={#1},breakable}
\newtcolorbox{saida}[1][]{colback=black!88,colframe=deitelblue,
  fontupper=\ttfamily\small\color{green!70},
  title={\small\color{white}Saída do programa},breakable,#1}

\setlist[enumerate,1]{label=\textbf{\alph*)},leftmargin=2em}
\setlist[itemize]{leftmargin=2em}

\lstset{
  basicstyle=\ttfamily\small,backgroundcolor=\color{deitelgray},
  frame=single,framerule=0.4pt,rulecolor=\color{deitelblue!40},
  breaklines=true,showstringspaces=false,
  keywordstyle=\color{deitelblue}\bfseries,
  commentstyle=\color{deitelgreen}\itshape,
  stringstyle=\color{deitelred},
  language=C,numbers=left,numberstyle=\tiny\color{gray},
  xleftmargin=18pt,xrightmargin=5pt,stepnumber=1,
}

\pagestyle{fancy}\fancyhf{}
\fancyhead[L]{\small\color{deitelblue}\textbf{C: Como Programar -- 6ª ed. | Deitel \& Deitel}}
\fancyhead[R]{\small\color{deitelblue}Cap. 5 -- Exercícios}
\fancyfoot[C]{\small\thepage}
\renewcommand{\headrulewidth}{0.4pt}

\titleformat{\section}[block]{\large\bfseries\color{deitelblue}}{\thesection.}{0.5em}{}[\titlerule]
\titleformat{\subsection}[block]{\normalsize\bfseries\color{deitelblue!80}}{}{}{}

\hypersetup{colorlinks=true,linkcolor=deitelblue}

\begin{document}

\begin{titlepage}
  \centering\vspace*{2cm}
  {\color{deitelblue}\rule{\linewidth}{2pt}}\\[0.4cm]
  {\huge\bfseries\color{deitelblue}C: Como Programar}\\[0.2cm]
  {\large\color{deitelblue}6ª Edição $\cdot$ Paul Deitel \& Harvey Deitel}\\[0.2cm]
  {\color{deitelblue}\rule{\linewidth}{2pt}}\\[1cm]
  {\LARGE\bfseries Capítulo 5}\\[0.3cm]
  {\Large\color{deitelblue!80}Funções em C}\\[0.8cm]
  \begin{tcolorbox}[colback=deitellightblue,colframe=deitelblue,width=0.85\linewidth]
    \centering{\large\bfseries Exercícios (5.8--5.48)}\\[0.2cm]
    {\normalsize Resolvidos passo a passo, com código completo}
  \end{tcolorbox}
\end{titlepage}

\tableofcontents\newpage

% ============================================================
\section{Exercício 5.8 -- Valores de \texttt{x}}
% ============================================================

\begin{respostabox}
\begin{enumerate}[label=\textbf{\alph*)}]
\item \texttt{fabs(7.5)} $\to$ $|7{,}5|$ $=$ \textbf{7,5}
\item \texttt{floor(7.5)} $\to$ maior inteiro $\leq 7{,}5$ $=$ \textbf{7,0}
\item \texttt{fabs(0.0)} $\to$ \textbf{0,0}
\item \texttt{ceil(0.0)} $\to$ \textbf{0,0}
\item \texttt{fabs(-6.4)} $\to$ \textbf{6,4}
\item \texttt{ceil(-6.4)} $\to$ menor inteiro $\geq -6{,}4$ $=$ \textbf{-6,0}
\item \texttt{ceil( -fabs( -8 + floor(-5.5) ) )}\\
  Avaliando de dentro para fora:
  \begin{itemize}[noitemsep]
    \item \texttt{floor(-5.5)} $=$ -6
    \item \texttt{-8 + (-6)} $=$ -14
    \item \texttt{fabs(-14)} $=$ 14
    \item \texttt{-(14)} $=$ -14
    \item \texttt{ceil(-14)} $=$ \textbf{-14,0}
  \end{itemize}
\end{enumerate}

\begin{boaprática}
  Atenção a \texttt{floor} de números negativos: \texttt{floor(-5.5)} é \texttt{-6}
  (mais negativo), não \texttt{-5}. Já \texttt{ceil(-5.5)} é \texttt{-5}.
  Pense em \textit{piso} como ``arredondar para baixo na reta numérica''.
\end{boaprática}
\end{respostabox}

% ============================================================
\section{Exercício 5.9 -- Tarifas de Estacionamento}
% ============================================================

\begin{respostabox}
\textbf{Regras:}
\begin{itemize}[noitemsep]
  \item Tarifa mínima: R\$ 2,00 (até 3 horas)
  \item Acima de 3 horas: R\$ 0,50 adicional por hora ou fração
  \item Máximo: R\$ 10,00 em qualquer período de 24 horas
\end{itemize}

\begin{lstlisting}
/* Ex5.9: tarifa_estacionamento.c
   Calcula taxas de estacionamento para 3 clientes */
#include <stdio.h>
#include <math.h>

double calcularTaxas( double horas ); /* prototipo */

int main( void )
{
    double horas;          /* horas estacionado */
    double taxa;           /* tarifa para o cliente atual */
    double totalHoras = 0.0;
    double totalReceita = 0.0;
    int cliente;

    printf( "%5s%8s%10s\n", "Carro", "Horas", "Taxa" );

    for ( cliente = 1; cliente <= 3; ++cliente ) {
        printf( "Digite horas do carro %d: ", cliente );
        scanf( "%lf", &horas );

        taxa = calcularTaxas( horas );

        printf( "%5d%8.1f%10.2f\n", cliente, horas, taxa );
        totalHoras   += horas;
        totalReceita += taxa;
    }

    printf( "%5s%8.1f%10.2f\n", "TOTAL", totalHoras, totalReceita );
    return 0;
} /* fim da funcao main */

/* Calcula a tarifa segundo as regras */
double calcularTaxas( double horas )
{
    double taxa = 2.00;  /* taxa minima ate 3 horas */

    if ( horas > 3.0 ) {
        /* Cobra 0,50 por hora ou fracao acima de 3 horas */
        /* ceil arredonda fracao para a hora cheia seguinte */
        taxa += 0.50 * ceil( horas - 3.0 );
    }

    if ( taxa > 10.00 ) {
        taxa = 10.00;     /* limite maximo diario */
    }

    return taxa;
} /* fim da funcao calcularTaxas */
\end{lstlisting}

\textbf{Análise:} Para 1{,}5 horas, está dentro das 3h iniciais $\to$ R\$ 2,00.
Para 4 horas, $\text{ceil}(4{-}3) = 1$ hora extra $\to$ 2 + 0{,}50 = R\$ 2,50.
Para 24 horas, calculado seria $2 + 0{,}50 \times 21 = 12{,}50$, mas limita-se em
R\$ 10,00.
\end{respostabox}

\newpage

% ============================================================
\section{Exercício 5.10 -- Arredondando para Inteiro}
% ============================================================

\begin{respostabox}
\textbf{Truque:} \texttt{floor(x + 0.5)} arredonda \texttt{x} para o inteiro mais
próximo. Por quê? Se a parte fracionária for $\geq 0{,}5$, somar 0{,}5 ``empurra''
o número para o próximo inteiro.

\begin{lstlisting}
/* Ex5.10: arredonda.c */
#include <stdio.h>
#include <math.h>

int main( void )
{
    double x, arredondado;

    printf( "Digite numero (negativo para terminar): " );
    scanf( "%lf", &x );

    while ( x >= 0.0 ) {
        arredondado = floor( x + 0.5 );
        printf( "Original: %.4f -> Arredondado: %.0f\n",
                x, arredondado );

        printf( "Proximo (negativo para terminar): " );
        scanf( "%lf", &x );
    }

    return 0;
}
\end{lstlisting}
\end{respostabox}

% ============================================================
\section{Exercício 5.11 -- Arredondamento por Casa Decimal}
% ============================================================

\begin{respostabox}
\begin{lstlisting}
/* Ex5.11: arredonda_casas.c */
#include <stdio.h>
#include <math.h>

double arredInteiro    ( double x );
double arredDecimos    ( double x );
double arredCentesimos ( double x );
double arredMilesimos  ( double x );

int main( void )
{
    double x;
    printf( "Digite um numero: " );
    scanf( "%lf", &x );

    printf( "Original:    %.6f\n", x );
    printf( "Inteiro:     %.0f\n", arredInteiro( x ) );
    printf( "Decimos:     %.1f\n", arredDecimos( x ) );
    printf( "Centesimos:  %.2f\n", arredCentesimos( x ) );
    printf( "Milesimos:   %.3f\n", arredMilesimos( x ) );

    return 0;
}

double arredInteiro( double x )
    { return floor( x + 0.5 ); }

double arredDecimos( double x )
    { return floor( x * 10.0 + 0.5 ) / 10.0; }

double arredCentesimos( double x )
    { return floor( x * 100.0 + 0.5 ) / 100.0; }

double arredMilesimos( double x )
    { return floor( x * 1000.0 + 0.5 ) / 1000.0; }
\end{lstlisting}

\textbf{Lógica:} para arredondar a $n$ casas decimais, multiplique por $10^n$,
arredonde para inteiro, e divida de volta por $10^n$. Por exemplo, para arredondar
3{,}14159 a 2 casas: $3{,}14159 \times 100 = 314{,}159$; \texttt{floor(314{,}659)} = 314;
$314 / 100 = 3{,}14$.
\end{respostabox}

\newpage

% ============================================================
\section{Exercício 5.12 -- Perguntas sobre Aleatoriedade}
% ============================================================

\begin{respostabox}
\begin{enumerate}[label=\textbf{\alph*)}]

\item \textbf{Escolher números ``aleatoriamente''} significa selecioná-los de tal
modo que cada valor possível tenha igual probabilidade de ser escolhido (distribuição
uniforme), e que escolhas sucessivas sejam estatisticamente independentes umas das
outras. Computadores não geram aleatoriedade verdadeira -- usam algoritmos
determinísticos para produzir sequências que \textit{parecem} aleatórias
(pseudoaleatoriedade).

\item \texttt{rand} é útil para simular jogos de azar porque modela situações como
lançamento de dados, baralho, moeda -- em que o resultado depende do acaso. Sem
\texttt{rand}, todas as execuções produziriam o mesmo resultado, perdendo a
imprevisibilidade característica do jogo.

\item \textbf{Por que usar \texttt{srand}?} Sem \texttt{srand}, o gerador
\texttt{rand} produz sempre a \textbf{mesma sequência} a cada execução do programa.
\texttt{srand(time(NULL))} fornece uma semente diferente a cada vez (baseada no
relógio do sistema), garantindo sequências distintas.

\textbf{Quando \textit{não} randomizar?} Durante o desenvolvimento e depuração,
queremos resultados \textit{reproduzíveis} -- o programa deve se comportar
exatamente igual a cada execução. Nesse contexto, omitir \texttt{srand} (ou usar
uma semente fixa) é desejável para identificar e corrigir bugs.

\item \textbf{Escala e deslocamento:} \texttt{rand()} retorna valores entre 0 e
\texttt{RAND\_MAX}. Para gerar valores em outro intervalo $[a, b]$:
\begin{itemize}[noitemsep]
  \item \textbf{Escala:} restringir a faixa a $b - a + 1$ valores: \texttt{rand() \% (b - a + 1)}
  \item \textbf{Deslocamento:} mover para começar em $a$: \texttt{a + ...}
\end{itemize}
A fórmula completa: \texttt{a + rand() \% (b - a + 1)}.

\item \textbf{Por que simulações são úteis?} Permitem investigar fenômenos do mundo
real:
\begin{itemize}[noitemsep]
  \item Quando experimentos reais são caros (testes de carros, aviões)
  \item Perigosos (reações nucleares, epidemiologia)
  \item Lentos (evolução, mudança climática)
  \item Ou impossíveis (cenários ``e se?'')
\end{itemize}
A simulação permite testar hipóteses, prever resultados e tomar decisões
informadas em segundos, com custo mínimo.
\end{enumerate}
\end{respostabox}

% ============================================================
\section{Exercício 5.13 -- Geração de aleatórios em intervalos}
% ============================================================

\begin{respostabox}
Fórmula geral: \texttt{n = a + rand() \% (b - a + 1)}, gerando valores em $[a, b]$.

\begin{enumerate}[label=\textbf{\alph*)}]
\item $1 \leq n \leq 2$: \texttt{n = 1 + rand() \% 2;}
\item $1 \leq n \leq 100$: \texttt{n = 1 + rand() \% 100;}
\item $0 \leq n \leq 9$: \texttt{n = rand() \% 10;}
\item $1000 \leq n \leq 1112$: \texttt{n = 1000 + rand() \% 113;} \quad ($1112-1000+1 = 113$)
\item $-1 \leq n \leq 1$: \texttt{n = -1 + rand() \% 3;} \quad ($1-(-1)+1 = 3$)
\item $-3 \leq n \leq 11$: \texttt{n = -3 + rand() \% 15;} \quad ($11-(-3)+1 = 15$)
\end{enumerate}
\end{respostabox}

% ============================================================
\section{Exercício 5.14 -- Aleatórios de conjuntos específicos}
% ============================================================

\begin{respostabox}
\textbf{Estratégia:} cada conjunto é uma progressão aritmética (PA) com 5 termos.
Geramos um inteiro aleatório no intervalo apropriado e aplicamos a fórmula
$\text{primeiro\_termo} + i \cdot \text{razão}$.

\begin{enumerate}[label=\textbf{\alph*)}]
\item \textbf{2, 4, 6, 8, 10:} primeiro 2, razão 2, 5 termos:
\begin{lstlisting}
printf( "%d", 2 + 2 * ( rand() % 5 ) );
\end{lstlisting}

\item \textbf{3, 5, 7, 9, 11:} primeiro 3, razão 2, 5 termos:
\begin{lstlisting}
printf( "%d", 3 + 2 * ( rand() % 5 ) );
\end{lstlisting}

\item \textbf{6, 10, 14, 18, 22:} primeiro 6, razão 4, 5 termos:
\begin{lstlisting}
printf( "%d", 6 + 4 * ( rand() % 5 ) );
\end{lstlisting}
\end{enumerate}
\end{respostabox}

\newpage

% ============================================================
\section{Exercício 5.15 -- Cálculo de Hipotenusa}
% ============================================================

\begin{respostabox}
\begin{lstlisting}
/* Ex5.15: hipotenusa.c
   Calcula hipotenusa de triangulo retangulo */
#include <stdio.h>
#include <math.h>

double hipotenusa( double cateto1, double cateto2 );

int main( void )
{
    double s1, s2;
    int triangulo;

    /* Tabela 5.18: 3 triangulos a testar */
    double lado1[] = { 3.0, 5.0, 8.0 };
    double lado2[] = { 4.0, 12.0, 15.0 };

    printf( "%-10s%-10s%-10s%-15s\n",
            "Triangulo", "Lado 1", "Lado 2", "Hipotenusa" );

    for ( triangulo = 0; triangulo < 3; ++triangulo ) {
        s1 = lado1[ triangulo ];
        s2 = lado2[ triangulo ];
        printf( "%-10d%-10.1f%-10.1f%-15.2f\n",
                triangulo + 1, s1, s2, hipotenusa( s1, s2 ) );
    }

    return 0;
}

/* Calcula hipotenusa: sqrt(a^2 + b^2) */
double hipotenusa( double cateto1, double cateto2 )
{
    return sqrt( cateto1 * cateto1 + cateto2 * cateto2 );
}
\end{lstlisting}

\textbf{Resultado esperado:} (3,4) $\to$ 5,00; (5,12) $\to$ 13,00; (8,15) $\to$ 17,00.
Todas são triplas de Pitágoras (já vistas no Ex.~4.27).

\textit{Nota:} usamos arrays (\texttt{double lado1[] = \{...\}}) -- conceito que
veremos formalmente no Capítulo~6. Aqui, é apenas uma forma compacta de listar
os valores de teste.
\end{respostabox}

% ============================================================
\section{Exercício 5.16 -- Função \texttt{potenciaInt}}
% ============================================================

\begin{respostabox}
\begin{lstlisting}
/* Ex5.16: potencia_int.c
   Calcula base^expoente sem usar funcoes de <math.h> */
#include <stdio.h>

int potenciaInt( int base, int expoente );

int main( void )
{
    int b, e;
    printf( "Digite base e expoente positivo: " );
    scanf( "%d%d", &b, &e );
    printf( "%d ^ %d = %d\n", b, e, potenciaInt( b, e ) );
    return 0;
}

/* Calcula base elevado a expoente usando for */
int potenciaInt( int base, int expoente )
{
    int resultado = 1;
    int i;

    for ( i = 1; i <= expoente; ++i ) {
        resultado *= base;
    }

    return resultado;
}
\end{lstlisting}

\textbf{Verificação:} \texttt{potenciaInt(3, 4)} $\to$ $3 \times 3 \times 3 \times 3 = 81$.
\end{respostabox}

\newpage

% ============================================================
\section{Exercício 5.17 -- Função \texttt{multiple}}
% ============================================================

\begin{respostabox}
\begin{lstlisting}
/* Ex5.17: multiplo.c
   Verifica se segundo inteiro e multiplo do primeiro */
#include <stdio.h>

int multiple( int a, int b );

int main( void )
{
    int x, y;

    printf( "Digite dois inteiros (-1 -1 para terminar): " );
    scanf( "%d%d", &x, &y );

    while ( x != -1 || y != -1 ) {
        if ( multiple( x, y ) ) {
            printf( "%d e multiplo de %d\n", y, x );
        }
        else {
            printf( "%d NAO e multiplo de %d\n", y, x );
        }
        printf( "Proximos dois (-1 -1 para terminar): " );
        scanf( "%d%d", &x, &y );
    }

    return 0;
}

/* Retorna 1 se b e multiplo de a, senao 0 */
int multiple( int a, int b )
{
    if ( a == 0 ) {
        return 0;          /* divisao por zero */
    }
    return ( b % a == 0 );  /* avalia para 1 ou 0 */
}
\end{lstlisting}
\end{respostabox}

% ============================================================
\section{Exercício 5.18 -- Função \texttt{even}}
% ============================================================

\begin{respostabox}
\begin{lstlisting}
/* Ex5.18: par.c
   Determina se inteiro e par */
#include <stdio.h>

int even( int n );

int main( void )
{
    int x;
    printf( "Digite inteiro (0 para terminar): " );
    scanf( "%d", &x );

    while ( x != 0 ) {
        if ( even( x ) ) {
            printf( "%d e PAR\n", x );
        }
        else {
            printf( "%d e IMPAR\n", x );
        }
        printf( "Proximo (0 para terminar): " );
        scanf( "%d", &x );
    }

    return 0;
}

int even( int n )
{
    return ( n % 2 == 0 );  /* 1 se par, 0 se impar */
}
\end{lstlisting}
\end{respostabox}

% ============================================================
\section{Exercício 5.19 -- Quadrado Sólido de Asteriscos}
% ============================================================

\begin{respostabox}
\begin{lstlisting}
/* Ex5.19: quadrado_asteriscos.c */
#include <stdio.h>

void exibeQuadrado( int lado );

int main( void )
{
    int n;
    printf( "Digite o lado: " );
    scanf( "%d", &n );

    exibeQuadrado( n );
    return 0;
}

void exibeQuadrado( int lado )
{
    int i, j;
    for ( i = 1; i <= lado; ++i ) {
        for ( j = 1; j <= lado; ++j ) {
            printf( "*" );
        }
        printf( "\n" );
    }
}
\end{lstlisting}
\end{respostabox}

% ============================================================
\section{Exercício 5.20 -- Quadrado de Caractere Variável}
% ============================================================

\begin{respostabox}
\begin{lstlisting}
/* Ex5.20: quadrado_caractere.c */
#include <stdio.h>

void exibeQuadrado( int lado, char fillCharacter );

int main( void )
{
    int n;
    char ch;

    printf( "Digite o lado e o caractere: " );
    scanf( "%d %c", &n, &ch );  /* espaco antes de %c */

    exibeQuadrado( n, ch );
    return 0;
}

void exibeQuadrado( int lado, char fillCharacter )
{
    int i, j;
    for ( i = 1; i <= lado; ++i ) {
        for ( j = 1; j <= lado; ++j ) {
            printf( "%c", fillCharacter );
        }
        printf( "\n" );
    }
}
\end{lstlisting}
\end{respostabox}

\newpage

% ============================================================
\section{Exercício 5.21 -- Projeto: Desenhando Formas}
% ============================================================

\begin{respostabox}
\begin{lstlisting}
/* Ex5.21: formas.c
   Programa que desenha varias formas usando funcoes */
#include <stdio.h>

void quadrado    ( int lado, char ch );
void retangulo   ( int largura, int altura, char ch );
void triangulo   ( int altura, char ch );
void linha       ( int comprimento, char ch );

int main( void )
{
    int opcao, n;
    char ch;

    printf( "Menu de formas:\n" );
    printf( "  1 - Quadrado\n" );
    printf( "  2 - Retangulo\n" );
    printf( "  3 - Triangulo\n" );
    printf( "  4 - Linha\n" );
    printf( "Escolha: " );
    scanf( "%d", &opcao );

    switch ( opcao ) {
        case 1:
            printf( "Lado e caractere: " );
            scanf( "%d %c", &n, &ch );
            quadrado( n, ch );
            break;
        case 2: {
            int alt;
            printf( "Largura, altura, caractere: " );
            scanf( "%d %d %c", &n, &alt, &ch );
            retangulo( n, alt, ch );
            break;
        }
        case 3:
            printf( "Altura e caractere: " );
            scanf( "%d %c", &n, &ch );
            triangulo( n, ch );
            break;
        case 4:
            printf( "Comprimento e caractere: " );
            scanf( "%d %c", &n, &ch );
            linha( n, ch );
            break;
        default:
            printf( "Opcao invalida\n" );
    }

    return 0;
}

void quadrado( int lado, char ch )
{
    int i, j;
    for ( i = 1; i <= lado; ++i ) {
        for ( j = 1; j <= lado; ++j ) printf( "%c", ch );
        printf( "\n" );
    }
}

void retangulo( int largura, int altura, char ch )
{
    int i, j;
    for ( i = 1; i <= altura; ++i ) {
        for ( j = 1; j <= largura; ++j ) printf( "%c", ch );
        printf( "\n" );
    }
}

void triangulo( int altura, char ch )
{
    int i, j;
    for ( i = 1; i <= altura; ++i ) {
        for ( j = 1; j <= i; ++j ) printf( "%c", ch );
        printf( "\n" );
    }
}

void linha( int comprimento, char ch )
{
    int j;
    for ( j = 1; j <= comprimento; ++j ) printf( "%c", ch );
    printf( "\n" );
}
\end{lstlisting}
\end{respostabox}

% ============================================================
\section{Exercício 5.22 -- Separando Dígitos}
% ============================================================

\begin{respostabox}
\begin{lstlisting}
/* Ex5.22: separa_digitos.c */
#include <stdio.h>

int parteInteira( int a, int b );  /* parte (a) */
int moduloInt   ( int a, int b );  /* parte (b) */
void exibeDigitos( int n );        /* parte (c) */

int main( void )
{
    int n;
    printf( "Digite inteiro entre 1 e 32767: " );
    scanf( "%d", &n );
    exibeDigitos( n );
    return 0;
}

int parteInteira( int a, int b )
{
    return a / b;
}

int moduloInt( int a, int b )
{
    return a % b;
}

/* Exibe digitos do numero separados por 2 espacos */
void exibeDigitos( int n )
{
    int divisor = 10000;  /* maior potencia <= 32767 */
    int digito;
    int comecou = 0;     /* para suprimir zeros a esquerda */

    while ( divisor > 0 ) {
        digito = parteInteira( n, divisor );
        n = moduloInt( n, divisor );

        if ( digito != 0 || comecou || divisor == 1 ) {
            printf( "%d  ", digito );
            comecou = 1;
        }

        divisor /= 10;
    }
    printf( "\n" );
}
\end{lstlisting}
\end{respostabox}

\newpage

% ============================================================
\section{Exercício 5.23 -- Tempo em Segundos}
% ============================================================

\begin{respostabox}
\begin{lstlisting}
/* Ex5.23: tempo_segundos.c
   Converte hora:minuto:segundo em segundos desde meia-noite */
#include <stdio.h>

int paraSegundos( int h, int m, int s );

int main( void )
{
    int h1, m1, s1, h2, m2, s2;
    int seg1, seg2, diferenca;

    printf( "Primeiro horario (h m s): " );
    scanf( "%d%d%d", &h1, &m1, &s1 );

    printf( "Segundo horario (h m s): " );
    scanf( "%d%d%d", &h2, &m2, &s2 );

    seg1 = paraSegundos( h1, m1, s1 );
    seg2 = paraSegundos( h2, m2, s2 );
    diferenca = seg2 - seg1;
    if ( diferenca < 0 ) diferenca = -diferenca;

    printf( "Diferenca: %d segundos\n", diferenca );
    return 0;
}

int paraSegundos( int h, int m, int s )
{
    return h * 3600 + m * 60 + s;
}
\end{lstlisting}
\end{respostabox}

% ============================================================
\section{Exercício 5.24 -- Conversões de Temperatura}
% ============================================================

\begin{respostabox}
\textbf{Fórmulas:}
\begin{align*}
  C &= \frac{5}{9}(F - 32) \\
  F &= \frac{9}{5}C + 32
\end{align*}

\begin{lstlisting}
/* Ex5.24: temperaturas.c */
#include <stdio.h>

int celsius( int fahrenheit );
int fahrenheit( int celsius );

int main( void )
{
    int t;

    printf( "Fahrenheit -> Celsius      Celsius -> Fahrenheit\n" );
    printf( "----------------------     ---------------------\n" );

    for ( t = 32; t <= 212; t += 18 ) {
        printf( "%4d F = %4d C            %4d C = %4d F\n",
                t, celsius( t ),
                t - 32, fahrenheit( t - 32 ) );
    }

    return 0;
}

int celsius( int f )
{
    return ( int )( 5.0 / 9.0 * ( f - 32 ) );
}

int fahrenheit( int c )
{
    return ( int )( 9.0 / 5.0 * c + 32 );
}
\end{lstlisting}

\textit{Nota:} usamos \textit{cast} para \texttt{int} para retornar o resultado
truncado, conforme exigido pelo enunciado (``com inteiros'').
\end{respostabox}

% ============================================================
\section{Exercício 5.25 -- Mínimo de 3 Números}
% ============================================================

\begin{respostabox}
\begin{lstlisting}
/* Ex5.25: minimo_tres.c */
#include <stdio.h>

double menor( double a, double b, double c );

int main( void )
{
    double x, y, z;
    printf( "Digite 3 numeros: " );
    scanf( "%lf%lf%lf", &x, &y, &z );

    printf( "O menor e: %f\n", menor( x, y, z ) );
    return 0;
}

double menor( double a, double b, double c )
{
    double m = a;
    if ( b < m ) m = b;
    if ( c < m ) m = c;
    return m;
}
\end{lstlisting}

\textbf{Versão idiomática com operador ternário:}
\begin{lstlisting}
double menor( double a, double b, double c )
{
    return ( a < b ? ( a < c ? a : c ) : ( b < c ? b : c ) );
}
\end{lstlisting}
\end{respostabox}

\newpage

% ============================================================
\section{Exercício 5.26 -- Números Perfeitos}
% ============================================================

\begin{respostabox}
\textbf{Definição:} Um inteiro $n$ é \textit{perfeito} se a soma de seus divisores
próprios (excluindo $n$) é igual a $n$.

\begin{lstlisting}
/* Ex5.26: numeros_perfeitos.c
   Encontra numeros perfeitos entre 1 e 1000 */
#include <stdio.h>

int perfect( int n );
void exibeFatores( int n );

int main( void )
{
    int n;

    printf( "Numeros perfeitos entre 1 e 1000:\n" );
    printf( "%-10s%s\n", "Numero", "Fatores" );

    for ( n = 2; n <= 1000; ++n ) {
        if ( perfect( n ) ) {
            printf( "%-10d", n );
            exibeFatores( n );
            printf( "\n" );
        }
    }

    return 0;
}

/* Retorna 1 se n e perfeito, senao 0 */
int perfect( int n )
{
    int soma = 1;        /* 1 e sempre fator (excluindo n) */
    int i;

    for ( i = 2; i * i <= n; ++i ) {
        if ( n % i == 0 ) {
            soma += i;
            if ( i != n / i ) {  /* evita contar raiz quadrada 2x */
                soma += n / i;
            }
        }
    }

    return ( n > 1 && soma == n );
}

void exibeFatores( int n )
{
    int i;
    printf( "1 " );
    for ( i = 2; i <= n / 2; ++i ) {
        if ( n % i == 0 ) printf( "%d ", i );
    }
}
\end{lstlisting}

\textbf{Resultado:} Os primeiros números perfeitos são \textbf{6} ($1+2+3$),
\textbf{28} ($1+2+4+7+14$), \textbf{496}, e \textbf{8128}.
\end{respostabox}

% ============================================================
\section{Exercício 5.27 -- Números Primos}
% ============================================================

\begin{respostabox}
\begin{lstlisting}
/* Ex5.27: primos.c */
#include <stdio.h>
#include <math.h>

int primo( int n );

int main( void )
{
    int n;
    int contador = 0;

    printf( "Primos de 1 a 10000:\n" );

    for ( n = 2; n <= 10000; ++n ) {
        if ( primo( n ) ) {
            ++contador;
            printf( "%6d", n );
            if ( contador % 10 == 0 ) printf( "\n" );
        }
    }

    printf( "\n\nTotal de primos: %d\n", contador );
    return 0;
}

/* Retorna 1 se n e primo, senao 0
   So precisa testar ate sqrt(n) -- ver explicacao abaixo */
int primo( int n )
{
    int i;
    int limite = ( int )sqrt( ( double )n );

    if ( n < 2 ) return 0;
    if ( n == 2 ) return 1;
    if ( n % 2 == 0 ) return 0;  /* pares > 2 nao sao primos */

    for ( i = 3; i <= limite; i += 2 ) {
        if ( n % i == 0 ) return 0;
    }

    return 1;
}
\end{lstlisting}

\textbf{Por que apenas até $\sqrt{n}$?} Se $n = a \times b$ e ambos $a, b > \sqrt{n}$,
então $a \times b > n$, contradição. Portanto, ao menos um dos fatores deve ser
$\leq \sqrt{n}$. Se não encontrarmos divisor até $\sqrt{n}$, $n$ é primo.

\textbf{Ganho de desempenho:} para $n = 10\,000$, $\sqrt{n} = 100$. Em vez de
testar 5\,000 valores ($n/2$), testamos apenas 100 -- redução de 50×!
\end{respostabox}

\newpage

% ============================================================
\section{Exercício 5.28 -- Inverter Dígitos}
% ============================================================

\begin{respostabox}
\begin{lstlisting}
/* Ex5.28: inverte_digitos.c */
#include <stdio.h>

int inverter( int n );

int main( void )
{
    int n;
    printf( "Digite um inteiro: " );
    scanf( "%d", &n );
    printf( "Invertido: %d\n", inverter( n ) );
    return 0;
}

int inverter( int n )
{
    int invertido = 0;
    while ( n > 0 ) {
        invertido = invertido * 10 + n % 10;
        n /= 10;
    }
    return invertido;
}
\end{lstlisting}

\textbf{Trace para 7631:} $i=0$, $n=7631$ $\to$ $i=1$, $n=763$ $\to$ $i=13$, $n=76$
$\to$ $i=136$, $n=7$ $\to$ $i=1367$, $n=0$. Resultado: \textbf{1367}.
\end{respostabox}

% ============================================================
\section{Exercício 5.29 -- MDC (Algoritmo de Euclides)}
% ============================================================

\begin{respostabox}
\textbf{Algoritmo de Euclides:} $\gcd(a, b) = \gcd(b, a \mod b)$ até que $b = 0$.

\begin{lstlisting}
/* Ex5.29: mdc.c */
#include <stdio.h>

int mdc( int a, int b );

int main( void )
{
    int x, y;
    printf( "Digite 2 inteiros positivos: " );
    scanf( "%d%d", &x, &y );
    printf( "MDC(%d, %d) = %d\n", x, y, mdc( x, y ) );
    return 0;
}

int mdc( int a, int b )
{
    int r;
    while ( b != 0 ) {
        r = a % b;
        a = b;
        b = r;
    }
    return a;
}
\end{lstlisting}

\textbf{Exemplo:} $\gcd(48, 36) = \gcd(36, 12) = \gcd(12, 0) = 12$.
\end{respostabox}

% ============================================================
\section{Exercício 5.30 -- Pontos de Qualidade}
% ============================================================

\begin{respostabox}
\begin{lstlisting}
/* Ex5.30: quality_points.c */
#include <stdio.h>

int qualityPoints( int media );

int main( void )
{
    int m;
    printf( "Digite a media (0-100): " );
    scanf( "%d", &m );
    printf( "Pontos de qualidade: %d\n", qualityPoints( m ) );
    return 0;
}

int qualityPoints( int media )
{
    if ( media >= 90 ) return 4;
    if ( media >= 80 ) return 3;
    if ( media >= 70 ) return 2;
    if ( media >= 60 ) return 1;
    return 0;
}
\end{lstlisting}
\end{respostabox}

\newpage

% ============================================================
\section{Exercício 5.31 -- Lançamento de Moeda}
% ============================================================

\begin{respostabox}
\begin{lstlisting}
/* Ex5.31: moeda.c
   Simula 100 lancamentos de moeda */
#include <stdio.h>
#include <stdlib.h>
#include <time.h>

int flip( void );

int main( void )
{
    int caras = 0;
    int coroas = 0;
    int i;

    srand( time( NULL ) );  /* randomiza */

    for ( i = 1; i <= 100; ++i ) {
        if ( flip() == 0 ) {
            ++caras;
            printf( "Cara  " );
        }
        else {
            ++coroas;
            printf( "Coroa " );
        }
        if ( i % 5 == 0 ) printf( "\n" );
    }

    printf( "\nTotal: %d caras, %d coroas\n", caras, coroas );
    return 0;
}

int flip( void )
{
    return rand() % 2;  /* 0 ou 1 com igual probabilidade */
}
\end{lstlisting}

\textbf{Resultado esperado:} aproximadamente 50 caras e 50 coroas (variação típica
de $\pm 5$ a $\pm 10$ pela natureza pseudoaleatória).
\end{respostabox}

% ============================================================
\section{Exercício 5.32 -- Descubra o Número}
% ============================================================

\begin{respostabox}
\begin{lstlisting}
/* Ex5.32: descubra_numero.c */
#include <stdio.h>
#include <stdlib.h>
#include <time.h>

int main( void )
{
    int numeroSecreto;
    int tentativa;
    char resposta;

    srand( time( NULL ) );

    do {
        numeroSecreto = 1 + rand() % 1000;

        printf( "Eu tenho um numero entre 1 e 1000.\n" );
        printf( "Voce consegue descobrir qual e?\n" );
        printf( "Digite sua primeira tentativa: " );
        scanf( "%d", &tentativa );

        while ( tentativa != numeroSecreto ) {
            if ( tentativa < numeroSecreto ) {
                printf( "Muito baixo. Tente novamente.\n" );
            }
            else {
                printf( "Muito alto. Tente novamente.\n" );
            }
            scanf( "%d", &tentativa );
        }

        printf( "Excelente! Voce descobriu o numero!\n" );
        printf( "Gostaria de jogar novamente (s ou n)? " );
        scanf( " %c", &resposta );

    } while ( resposta == 's' || resposta == 'S' );

    return 0;
}
\end{lstlisting}
\end{respostabox}

\newpage

% ============================================================
\section{Exercício 5.33 -- Modificação de ``Descubra o Número''}
% ============================================================

\begin{respostabox}
\textbf{Por que 10 ou menos tentativas bastam?}

A técnica usada é \textbf{busca binária}: a cada tentativa, dividimos o intervalo
ao meio. Como $2^{10} = 1024 > 1000$, em \textbf{no máximo 10} divisões pelo meio
consegue-se localizar qualquer número de 1 a 1000.

\textit{Estratégia ótima:} chute sempre o meio do intervalo restante. Começa com
500. Se ``muito alto'', chute 250. Se ``muito baixo'', chute 125. E assim por diante.

\begin{lstlisting}
/* Ex5.33: modificacao do jogo */
#include <stdio.h>
#include <stdlib.h>
#include <time.h>

int main( void )
{
    int numeroSecreto, tentativa, contador;
    char resposta;

    srand( time( NULL ) );

    do {
        numeroSecreto = 1 + rand() % 1000;
        contador = 0;

        printf( "Tente descobrir o numero (1 a 1000): " );

        do {
            scanf( "%d", &tentativa );
            ++contador;

            if ( tentativa < numeroSecreto ) {
                printf( "Muito baixo. Tente novamente: " );
            }
            else if ( tentativa > numeroSecreto ) {
                printf( "Muito alto. Tente novamente: " );
            }
        } while ( tentativa != numeroSecreto );

        printf( "Excelente! Acertou em %d tentativas.\n", contador );

        if ( contador < 10 ) {
            printf( "Ou voce conhece o segredo ou teve sorte!\n" );
        }
        else if ( contador == 10 ) {
            printf( "Ahah! Voce conhece o segredo!\n" );
        }
        else {
            printf( "Voce deveria se sair melhor!\n" );
        }

        printf( "Jogar novamente? (s/n): " );
        scanf( " %c", &resposta );
    } while ( resposta == 's' || resposta == 'S' );

    return 0;
}
\end{lstlisting}
\end{respostabox}

\newpage

% ============================================================
\section{Exercício 5.34 -- Exponenciação Recursiva}
% ============================================================

\begin{respostabox}
\textbf{Definição recursiva:} $\text{base}^{\text{exp}} = \text{base} \times \text{base}^{\text{exp}-1}$,
com caso básico $\text{base}^1 = \text{base}$.

\begin{lstlisting}
/* Ex5.34: power_recursiva.c */
#include <stdio.h>

int power( int base, int expoente );

int main( void )
{
    int b, e;
    printf( "Digite base e expoente: " );
    scanf( "%d%d", &b, &e );
    printf( "%d ^ %d = %d\n", b, e, power( b, e ) );
    return 0;
}

int power( int base, int expoente )
{
    if ( expoente == 1 ) {           /* caso basico */
        return base;
    }
    return base * power( base, expoente - 1 );  /* recursao */
}
\end{lstlisting}

\textbf{Trace para \texttt{power(3, 4)}:}
\begin{itemize}[noitemsep]
  \item \texttt{power(3,4)} = $3 \times $ \texttt{power(3,3)}
  \item \texttt{power(3,3)} = $3 \times $ \texttt{power(3,2)}
  \item \texttt{power(3,2)} = $3 \times $ \texttt{power(3,1)}
  \item \texttt{power(3,1)} = 3 (caso básico)
\end{itemize}
Resultado: $3 \times 3 \times 3 \times 3 = 81$.
\end{respostabox}

% ============================================================
\section{Exercício 5.35 -- Fibonacci Iterativo com \texttt{double}}
% ============================================================

\begin{respostabox}
\textbf{Versão não recursiva:}
\begin{lstlisting}
/* Ex5.35: fibonacci_iterativo.c */
#include <stdio.h>

double fibonacci( int n );

int main( void )
{
    int n = 0;

    /* Loop ate o resultado falhar (overflow) */
    while ( 1 ) {
        printf( "fib(%d) = %.0f\n", n, fibonacci( n ) );
        ++n;
    }

    return 0;
}

double fibonacci( int n )
{
    double anterior = 0.0;
    double atual = 1.0;
    double proximo;
    int i;

    if ( n == 0 ) return 0.0;
    if ( n == 1 ) return 1.0;

    for ( i = 2; i <= n; ++i ) {
        proximo = anterior + atual;
        anterior = atual;
        atual = proximo;
    }

    return atual;
}
\end{lstlisting}

\textbf{Limites:} \texttt{double} representa até cerca de $1{,}8 \times 10^{308}$.
$F_{1476} \approx 1{,}3 \times 10^{308}$, então em torno desse índice começam os
\textit{overflows} (resultando em \texttt{inf}).
\end{respostabox}

\newpage

% ============================================================
\section{Exercício 5.36 -- Torres de Hanói}
% ============================================================

\begin{respostabox}
\textbf{Decomposição recursiva:} Para mover $n$ discos do pino A para o pino C usando
B como auxiliar:
\begin{enumerate}[noitemsep]
  \item Mover $n-1$ discos de A para B (usando C como auxiliar).
  \item Mover o maior disco de A para C.
  \item Mover $n-1$ discos de B para C (usando A como auxiliar).
\end{enumerate}

\begin{lstlisting}
/* Ex5.36: torres_hanoi.c */
#include <stdio.h>

void hanoi( int n, int origem, int destino, int auxiliar );

int main( void )
{
    int n;
    printf( "Numero de discos: " );
    scanf( "%d", &n );

    printf( "Sequencia de movimentos:\n" );
    hanoi( n, 1, 3, 2 );  /* origem=1, destino=3, auxiliar=2 */

    return 0;
}

void hanoi( int n, int origem, int destino, int auxiliar )
{
    if ( n == 1 ) {                                   /* caso basico */
        printf( "%d -> %d\n", origem, destino );
        return;
    }
    /* Move n-1 discos de origem para auxiliar */
    hanoi( n - 1, origem, auxiliar, destino );

    /* Move o disco maior */
    printf( "%d -> %d\n", origem, destino );

    /* Move n-1 discos de auxiliar para destino */
    hanoi( n - 1, auxiliar, destino, origem );
}
\end{lstlisting}

\textbf{Saída para n=3:}
\begin{saida}
1 -> 3\\1 -> 2\\3 -> 2\\1 -> 3\\2 -> 1\\2 -> 3\\1 -> 3
\end{saida}

\textbf{Curiosidade:} para $n$ discos são necessários $2^n - 1$ movimentos. Para
os 64 discos da lenda, $2^{64} - 1 \approx 1{,}8 \times 10^{19}$ movimentos --
mesmo a 1 movimento por segundo, levaria $\sim 585$ bilhões de anos!
\end{respostabox}

% ============================================================
\section{Exercício 5.38 -- Visualizando a Recursão (Fatorial)}
% ============================================================

\begin{respostabox}
\begin{lstlisting}
/* Ex5.38: fatorial_visual.c */
#include <stdio.h>

long fatorial( int n, int nivel );

int main( void )
{
    int n;
    printf( "Digite n: " );
    scanf( "%d", &n );
    printf( "Resultado: %ld\n", fatorial( n, 0 ) );
    return 0;
}

long fatorial( int n, int nivel )
{
    int i;
    long resultado;

    /* Recuo proporcional ao nivel */
    for ( i = 0; i < nivel; ++i ) printf( "  " );
    printf( "fatorial(%d) chamado\n", n );

    if ( n <= 1 ) {
        for ( i = 0; i < nivel; ++i ) printf( "  " );
        printf( "fatorial(%d) retorna 1\n", n );
        return 1;
    }

    resultado = n * fatorial( n - 1, nivel + 1 );

    for ( i = 0; i < nivel; ++i ) printf( "  " );
    printf( "fatorial(%d) retorna %ld\n", n, resultado );

    return resultado;
}
\end{lstlisting}

\textbf{Saída para n=4:}
\begin{saida}
fatorial(4) chamado\\
\hphantom{xx}fatorial(3) chamado\\
\hphantom{xxxx}fatorial(2) chamado\\
\hphantom{xxxxxx}fatorial(1) chamado\\
\hphantom{xxxxxx}fatorial(1) retorna 1\\
\hphantom{xxxx}fatorial(2) retorna 2\\
\hphantom{xx}fatorial(3) retorna 6\\
fatorial(4) retorna 24
\end{saida}
\end{respostabox}

\newpage

% ============================================================
\section{Exercício 5.39 -- MDC Recursivo}
% ============================================================

\begin{respostabox}
\begin{lstlisting}
/* Ex5.39: mdc_recursivo.c */
#include <stdio.h>

int mdc( int x, int y );

int main( void )
{
    int a, b;
    printf( "Digite 2 inteiros: " );
    scanf( "%d%d", &a, &b );
    printf( "MDC(%d, %d) = %d\n", a, b, mdc( a, b ) );
    return 0;
}

int mdc( int x, int y )
{
    if ( y == 0 ) return x;        /* caso basico */
    return mdc( y, x % y );        /* recursao */
}
\end{lstlisting}

\textbf{Trace para \texttt{mdc(48, 36)}:}
\begin{itemize}[noitemsep]
  \item \texttt{mdc(48, 36)} $\to$ \texttt{mdc(36, 48 \% 36)} = \texttt{mdc(36, 12)}
  \item \texttt{mdc(36, 12)} $\to$ \texttt{mdc(12, 0)}
  \item \texttt{mdc(12, 0)} $\to$ retorna 12 (caso básico)
\end{itemize}
Resultado: \textbf{12}.
\end{respostabox}

% ============================================================
\section{Exercício 5.40 -- \texttt{main} Recursiva}
% ============================================================

\begin{respostabox}
\begin{lstlisting}
/* Ex5.40: main_recursiva.c */
#include <stdio.h>

int main( void )
{
    static int count = 1;
    printf( "Chamada #%d\n", count );
    count++;
    main();           /* chama-se recursivamente */
    return 0;
}
\end{lstlisting}

\textbf{O que acontece?} A função \texttt{main} \textbf{pode} chamar a si mesma em
C (não há restrição da linguagem para isso). O programa imprime sucessivamente:
\texttt{Chamada \#1}, \texttt{Chamada \#2}, \texttt{Chamada \#3}\ldots até que a
\textbf{pilha de execução transborde} (\textit{stack overflow}). Cada chamada
consome espaço de pilha; eventualmente, o sistema operacional encerra o programa
com falha de segmentação.

\textit{Nota:} embora compile e execute, é uma prática \textbf{altamente desencorajada}.
A recursão deve ter um caso básico que termine; aqui não há.
\end{respostabox}

% ============================================================
\section{Exercício 5.41 -- Distância entre Pontos}
% ============================================================

\begin{respostabox}
\begin{lstlisting}
/* Ex5.41: distancia.c */
#include <stdio.h>
#include <math.h>

double distance( double x1, double y1, double x2, double y2 );

int main( void )
{
    double x1, y1, x2, y2;
    printf( "Digite o primeiro ponto (x1 y1): " );
    scanf( "%lf%lf", &x1, &y1 );
    printf( "Digite o segundo ponto (x2 y2): " );
    scanf( "%lf%lf", &x2, &y2 );
    printf( "Distancia: %.4f\n",
            distance( x1, y1, x2, y2 ) );
    return 0;
}

double distance( double x1, double y1, double x2, double y2 )
{
    double dx = x2 - x1;
    double dy = y2 - y1;
    return sqrt( dx * dx + dy * dy );
}
\end{lstlisting}
\end{respostabox}

\newpage

% ============================================================
\section{Exercício 5.42 -- O que o programa faz?}
% ============================================================

\begin{respostabox}
\begin{lstlisting}
if ( ( c = getchar() ) != EOF ) {
    main();
    printf( "%c", c );
}
\end{lstlisting}

\textbf{Análise:} O programa lê caracteres do teclado um a um e usa recursão
através de \texttt{main()}. Para cada caractere lido, antes de imprimi-lo, chama
\texttt{main()} recursivamente para ler e processar o próximo. Quando \texttt{EOF}
é encontrado, a recursão começa a desempilhar -- imprimindo os caracteres
\textbf{na ordem inversa}.

\textbf{Conclusão:} o programa \textbf{inverte uma sequência de caracteres da entrada}.

\textbf{Exemplo:} entrada \texttt{ABC<EOF>} produz saída \texttt{CBA}.
\end{respostabox}

% ============================================================
\section{Exercício 5.43 -- O que faz \texttt{mystery}?}
% ============================================================

\begin{respostabox}
\begin{lstlisting}
int mystery( int a, int b )
{
    if ( b == 1 ) return a;
    else return a + mystery( a, b - 1 );
}
\end{lstlisting}

\textbf{Análise:} Para \texttt{mystery(a, b)}:
\begin{itemize}[noitemsep]
  \item Se $b = 1$: retorna $a$.
  \item Senão: retorna $a + \text{mystery}(a, b-1)$.
\end{itemize}

Expandindo: \texttt{mystery(a, b)} = $a + a + \ldots + a$ ($b$ vezes) = $a \times b$.

\textbf{Conclusão:} a função calcula a \textbf{multiplicação $a \times b$ por adições
sucessivas} -- uma forma recursiva clássica de implementar multiplicação.

\textbf{Restrição:} $b$ deve ser positivo para evitar recursão infinita.
\end{respostabox}

% ============================================================
\section{Exercício 5.44 -- \texttt{mystery} sem restrição em \texttt{b}}
% ============================================================

\begin{respostabox}
Para tratar $b$ negativo ou zero:

\begin{lstlisting}
int mystery( int a, int b )
{
    if ( b == 0 ) return 0;        /* a * 0 = 0 */
    if ( b < 0 ) {
        return -mystery( a, -b );  /* a * (-b) = -(a * b) */
    }
    if ( b == 1 ) return a;
    return a + mystery( a, b - 1 );
}
\end{lstlisting}

\textbf{Lógica:}
\begin{itemize}[noitemsep]
  \item Se $b = 0$, o produto é 0.
  \item Se $b < 0$, calculamos com $|b|$ e negamos o resultado.
  \item Caso contrário, mantemos a recursão original.
\end{itemize}
\end{respostabox}

\newpage

% ============================================================
\section{Exercício 5.46 -- Identificar e Corrigir Erros}
% ============================================================

\begin{respostabox}
\begin{enumerate}[label=\textbf{\alph*)}]

\item \textbf{Erro:} O cabeçalho da definição não tem o tipo de retorno -- deveria
ser \texttt{double cube(float number)}.\\
\textbf{Correção:}
\begin{lstlisting}
double cube( float ); /* prototipo */
double cube( float number ) /* definicao */
{
    return number * number * number;
}
\end{lstlisting}

\item \textbf{Erro:} Combinar dois especificadores de classe de armazenamento
(\texttt{register} e \texttt{auto}) na mesma declaração é proibido.\\
\textbf{Correção:} usar apenas um:
\begin{lstlisting}
register int x = 7;  /* OU: auto int x = 7; */
\end{lstlisting}

\item \textbf{Erro:} \texttt{srand} é uma função \texttt{void} (não retorna valor)
e exige um argumento. Não pode ser usada em uma atribuição como aqui.\\
\textbf{Correção:} provavelmente o autor quis usar \texttt{rand}:
\begin{lstlisting}
int randomNumber = rand();  /* obtem um aleatorio */
srand( time( NULL ) );      /* uso correto de srand */
\end{lstlisting}

\item \textbf{Erro:} A conversão \texttt{(double) x} no \texttt{printf} é redundante
e fora de propósito. O verdadeiro problema é exibir um \texttt{int} com \texttt{\%f},
o que é incorreto.\\
\textbf{Correção:}
\begin{lstlisting}
double y = 123.45678;
int x;
x = (int) y;       /* cast explicito recomendado */
printf( "%d\n", x );
\end{lstlisting}

\item \textbf{Erro:} A variável \texttt{number} é \textit{redeclarada} dentro da
função, sobrepondo o parâmetro do mesmo nome (Erro Comum 5.5).\\
\textbf{Correção:} remover a redeclaração:
\begin{lstlisting}
double square( double number )
{
    return number * number;  /* sem 'double number;' aqui */
}
\end{lstlisting}

\item \textbf{Erro:} Na recursão, a chamada \texttt{sum(n)} usa o mesmo \texttt{n},
sem reduzir -- gera \textbf{recursão infinita}.\\
\textbf{Correção:}
\begin{lstlisting}
int sum( int n )
{
    if ( n == 0 ) return 0;
    return n + sum( n - 1 );  /* deve reduzir o argumento */
}
\end{lstlisting}

\end{enumerate}
\end{respostabox}

% ============================================================
\section{Exercício 5.47 -- Modificação do Jogo Craps}
% ============================================================

\begin{respostabox}
\begin{lstlisting}
/* Ex5.47: craps_apostas.c
   Versao do jogo craps com sistema de apostas */
#include <stdio.h>
#include <stdlib.h>
#include <time.h>

int rolarDados( void );

enum Status { CONTINUA, VENCEU, PERDEU };

int main( void )
{
    int saldoBanca = 1000;
    int aposta;
    int soma, meuPonto;
    enum Status status;

    srand( time( NULL ) );

    while ( saldoBanca > 0 ) {
        printf( "\nSaldo: R$ %d\n", saldoBanca );
        printf( "Aposta (max %d): ", saldoBanca );
        scanf( "%d", &aposta );

        while ( aposta <= 0 || aposta > saldoBanca ) {
            printf( "Aposta invalida! Digite (max %d): ",
                    saldoBanca );
            scanf( "%d", &aposta );
        }

        soma = rolarDados();
        switch ( soma ) {
            case 7:
            case 11:
                status = VENCEU;
                break;
            case 2:
            case 3:
            case 12:
                status = PERDEU;
                break;
            default:
                status = CONTINUA;
                meuPonto = soma;
                printf( "Ponto e %d\n", meuPonto );
                break;
        }

        while ( status == CONTINUA ) {
            soma = rolarDados();
            if ( soma == meuPonto ) status = VENCEU;
            else if ( soma == 7 ) status = PERDEU;
        }

        if ( status == VENCEU ) {
            saldoBanca += aposta;
            printf( "Voce venceu! Novo saldo: R$ %d\n", saldoBanca );
            printf( "Voce esta ganhando. Aproveite suas fichas!\n" );
        }
        else {
            saldoBanca -= aposta;
            printf( "Voce perdeu. Novo saldo: R$ %d\n", saldoBanca );
            if ( saldoBanca == 0 ) {
                printf( "Sinto muito. Voce esta quebrado!\n" );
                break;
            }
            printf( "Vamos la, continue tentando!\n" );
        }
    }

    return 0;
}

int rolarDados( void )
{
    int dado1 = 1 + rand() % 6;
    int dado2 = 1 + rand() % 6;
    int soma = dado1 + dado2;
    printf( "Dados: %d + %d = %d\n", dado1, dado2, soma );
    return soma;
}
\end{lstlisting}
\end{respostabox}

\newpage

% ============================================================
\section{Exercício 5.48 -- Aperfeiçoando Fibonacci Recursivo}
% ============================================================

\begin{respostabox}
\textbf{O problema da implementação ingênua:}

A versão recursiva clássica:
\begin{lstlisting}
long fib( int n ) {
    if ( n <= 1 ) return n;
    return fib( n - 1 ) + fib( n - 2 );
}
\end{lstlisting}
tem complexidade \textbf{exponencial} ($O(\phi^n)$, onde $\phi \approx 1{,}618$),
porque \textit{recalcula} os mesmos valores muitas vezes. Para \texttt{fib(40)},
são feitas \(\approx 10^9\) chamadas!

\textbf{Técnicas de melhoria:}

\textbf{1. Versão iterativa (Ex.~5.35):}
\begin{lstlisting}
long fib( int n ) {
    long a = 0, b = 1, c;
    int i;
    for ( i = 2; i <= n; ++i ) {
        c = a + b;
        a = b;
        b = c;
    }
    return ( n == 0 ? 0 : b );
}
\end{lstlisting}
Complexidade $O(n)$, espaço $O(1)$.

\textbf{2. Recursão de cauda (\textit{tail recursion}):}
\begin{lstlisting}
long fibAux( int n, long a, long b ) {
    if ( n == 0 ) return a;
    return fibAux( n - 1, b, a + b );
}
long fib( int n ) { return fibAux( n, 0, 1 ); }
\end{lstlisting}
Compiladores podem otimizar (eliminar a chamada da pilha), tornando praticamente
equivalente à versão iterativa.

\textbf{3. Memoização:} armazenar resultados em uma tabela para evitar recálculos.
Reduz a complexidade da versão recursiva de $O(\phi^n)$ para $O(n)$, mas com custo
de memória $O(n)$.

\textbf{Comparação dos méritos:}
\begin{itemize}[noitemsep]
  \item \textbf{Iterativa:} mais rápida, menos memória, mais difícil de ler para
  problemas naturalmente recursivos.
  \item \textbf{Tail recursion:} elegância da recursão + eficiência da iteração
  (quando o compilador otimiza).
  \item \textbf{Memoização:} preserva clareza recursiva, mas exige memória
  adicional. Excelente para problemas com sobreposição de subproblemas.
\end{itemize}
\end{respostabox}

\end{document}
